\providecommand{\kBKM}{\kappa_{\mathrm{BKM}}}
\providecommand{\kch}{\kappa_{\mathrm{ch}}}
\providecommand{\kcat}{\kappa_{\mathrm{cat}}}
\providecommand{\cO}{\mathcal{O}}
\providecommand{\ZZ}{\mathbb{Z}}
\providecommand{\PP}{\mathbb{P}}
\providecommand{\gDelta}{\mathfrak{g}_{\Delta_5}}
\providecommand{\Spf}{\mathrm{Sp}_4(\ZZ)}

\section*{Slides lines 870--953: the eight-form enumeration}

\paragraph{What the slides state (lines 870--953).}
The final stretch enumerates $(S, g_N, h_M)$ data with $N, M \le 8$ and
constructs, for each $N$, a DT-zeta function $Z^X(q, t, p)$ of the
smooth projective CY3 $X = (S \times E)/(\ZZ/N\ZZ)$.  The DT series is
assembled as
\[
  Z^X(q, t, p) \;=\; \sum_{h \ge 0} \sum_{d \ge 0} \sum_{n \in \ZZ}
      \mathrm{DT}_{n, (\beta_h, d)}\; q^{d - 1} t^{\tfrac{1}{2}\langle \beta_h, \beta_h\rangle}(-p)^n,
\]
with classes $\beta_h = (s_1 + \cdots + s_N + h F)/N$ built from the
$N$ torsion sections $s_1, \ldots, s_N$ and the fibre class $F$ of
$\pi: S \to \PP^1$; Behrend-weighted Euler characteristics
$\mathrm{DT}_{n,(\beta_h,d)} = \int_{\mathrm{Hilb}^n(X,\beta_h)/E} \nu_{\mathrm{dE}}$
define the coefficients.  The theorem at $N = 1$ reads
\[
  Z^{S \times E} \;=\; -\,\frac{C}{(\Delta_5)^2},
  \qquad \text{K3 elliptic genus } = 2\phi_{0, 1},
\]
and the conjecture (slide line 929) advances this to all $N$: the
$L$-$(h_M = g_N^s)$-twisted DT series $-1/Z^X_{L, h_M}(q, t, p)$ is a
Siegel modular form, and there is a GKM algebra
$\mathfrak{g}_{(g_N, h_M)} \subset \mathfrak{g}_L$ with multiplicities
controlled by the $(g_N, h_M)$-twisted elliptic genus
$F_L^{(g_N, h_M)}$, a Jacobi form of weight $0$ and index $1$.

\paragraph{Final six slides: the case split.}
After defining $A = \tfrac{1}{4}\phi_{0,1}$ and
$B = \phi_{-2, 1} = \nu_1^2/\eta^6$ (line 939), the talk partitions
$N \in \{1, \ldots, 8\}$ into four case clusters and states the
twisted-genus formula in each:
\begin{enumerate}
\item[(a)] \emph{Generic cluster} $N \in \{1, 2, 3, 5, 7\}$
  (line 943): uniform formula valid for all $1 \le s, r \le N - 1$
  and $0 \le k \le N - 1$.  The slide frames this as the primary
  generic family — the same generic family where
  $g_N$ acts on $H^*(S, \ZZ)$ with spectrum determined by a single
  cyclotomic divisor pattern.
\item[(b)] $N = 4$ (line 945), enumerated over $s \in \{0, 1, 2, 3\}$.
\item[(c)] $N = 6$ (lines 947, 949), enumerated over $s \in \{0, 1, 2, 3\}$.
\item[(d)] $N = 8$ (lines 951, 953), without an $s$-enumeration on the
  slide header.
\end{enumerate}
The case split is \emph{arithmetic}, not physical: the primes
$\{2, 3, 5, 7\}$ together with $N = 1$ form one uniform family because
$g_N^r$ generates a subgroup of order $N/\gcd(r, N)$ whose character
sum decomposes through a single cyclotomic polynomial;
the composite orders $N \in \{4, 6, 8\}$ each require separate
orbit-sum formulas because multiple cyclotomic factors enter.

\paragraph{Cross-check against \textsc{Lorgat 2020}.}
The slide partition $\{1, 2, 3, 5, 7\} \cup \{4\} \cup \{6\} \cup \{8\}$
aligns with the \emph{eight Gritsenko--Cl\'ery paramodular diagonal-divisor
forms} — the subject of Conjecture 1 of Lorgat 2020
(\texttt{reference\_lorgat\_2020\_automorphic\_corrections.md}): the eight
paramodular forms indexed by $N \le 8$ are the eight denominator
functions the conjecture associates to $(S, g_N)$-twisted DT zetas,
one per case cluster plus the prime-cluster entries.
The slides state Conjecture 1 in prospective form (\emph{``all 8
Siegel modular forms arise as $-1/Z^X_{L, h_M}$''}, line 929), matching
exactly the paper's conjectural scope.

\paragraph{CHL-admissibility.}
\textbf{The slides contain no ``CHL-admissibility'' language.}  The
restriction to $N \in \{1, 2, 3, 4, 6\}$ that governs the Vol III
programme's Borcherds-weight identity
$\kBKM(\Phi_N) = c_N(0)/2$ is \emph{not} imposed in the slides: the
slides enumerate through all of $\{1, \ldots, 8\}$, and the
prime-cluster slide explicitly includes $N = 5, 7$.  This is a
non-trivial divergence of scope.

Reconciliation: the programme's $\{1, 2, 3, 4, 6\}$ set is the
\emph{CHL-admissible} subset — those $N$ for which $\ZZ/N\ZZ$ acts
\emph{symplectically} on a K3 lattice extending to a free action on
$S \times E$ with the quotient carrying the relevant
$\SU(2)$ R-symmetry structure needed for string-theoretic $1/2$-BPS
countings.  The slides work at the broader Gritsenko--Cl\'ery level
of $\Spf$ modular forms, where the constraint on $N$ is only
``$N \le 8$ and $g_N$ preserves the elliptic genus of $S$''.  The
slide cases $N \in \{5, 7\}$ and $N = 8$ correspond to paramodular
forms of the eight-form family that do \emph{not} admit the full
$(g_N, h_M)$ twist pairing satisfying the CHL physical consistency
condition; the Vol III programme treats these as outside the scope
of $\kBKM(\Phi_N) = c_N(0)/2$ (AP-CY on scope discipline).

\paragraph{In slides vs added by programme.}
\begin{itemize}
\item \emph{In slides} (lines 870--953): K3 elliptic fibration with
  $N$ sections, $(S \times E)/(\ZZ/N\ZZ)$ CY3 construction, DT zeta
  $Z^X(q, t, p)$ with explicit $\beta_h$, the $N = 1$ theorem
  $Z^{S \times E} = -C/(\Delta_5)^2$, Conjecture 1 enumeration across
  the four clusters $\{1, 2, 3, 5, 7\}, \{4\}, \{6\}, \{8\}$, and the
  $A = \phi_{0,1}/4$, $B = \phi_{-2,1}$ generators of the weak Jacobi
  ring $J_{0, *}$.
\item \emph{Added by programme, not in slides}:
  \begin{itemize}
  \item The identity $\kBKM(\Phi_N) = c_N(0)/2$ as a universal
    statement for $N \in \{1, 2, 3, 4, 6\}$ — slides state
    \emph{denominator identity} data but do not extract
    $\kappa$-invariants from weights.
  \item The CHL-admissibility cut $\{1, 2, 3, 4, 6\} \subset \{1, \ldots, 8\}$.
  \item The $\{2, 3, 5, 24\}$ constellation on $S \times E$: slides
    give only $\kcat(S \times E) = 0$ indirectly (via the K3 elliptic
    genus being $2\phi_{0,1}$), not the four-construction picture.
  \item The CY-D dimensional stratification.
  \item The cross-identification
    $\kBKM(\Phi_N) = \kch(K3 \times E) + \chi(\cO_{\text{fibre}})$, which
    is \emph{false} at every $N$ and is explicitly disclaimed in the
    cache; the slides never assert this identification.
  \end{itemize}
\end{itemize}

\paragraph{Summary (three lines).}
The final slide stretch inscribes the eight-$N$ case enumeration
underpinning Conjecture 1 of Lorgat 2020, with case clusters
$\{1, 2, 3, 5, 7\}, \{4\}, \{6\}, \{8\}$ covering the full
Gritsenko--Cl\'ery paramodular-form family.  The Vol III programme
\emph{subset} $\{1, 2, 3, 4, 6\}$ is a CHL-admissibility restriction
added by the programme, not present in the slide material.  All
$\kBKM$-weight arithmetic ($c_N(0)/2$ as the $\Phi_N$ weight) is
programme-internal and draws on slide content only through the
$N = 1$ theorem $Z^{S\times E} = -C/(\Delta_5)^2$ and the associated
$\phi_{0, 1}$ Fourier-coefficient data.
